\chapter{Results}
This chapter presents the results from the human perception study. The results are organized into 3 sections, each covering emotion/intent recognition, Likert-scale ratings, and Godspeed questionnaire results. Where relevant, the results of the live study (N=16) and video-based study (N=22) are presented separately and compared to identify any differences in perception between the two approaches.

\section{Emotion Recognition Results}
This section presents the results of the open-ended emotion recognition task, in which participants were asked to describe the emotion or intent the robot was conveying in each of the 6 expressive animations. As previously mentioned, the responses were divided into 3 groups: \textit{intended}, \textit{adjacent}, or \textit{other}. Figure~\ref{fig:recognition_rate} summarizes the recognition rate, i.e., the rate at which the responses were categorized as the intended emotion. IT shows the recognition rates across all six animations for both the live and video-based groups, followed by a detailed breakdown of the response distribution in figures~\ref{fig:emotion_recognition_live} and ~\ref{fig:emotion_recognition_video}.

%\newpage
\begin{figure}[H]
\centering
\begin{tikzpicture}
\begin{axis}[
    ybar,
    bar width=12pt,
    width=11cm,
    height=8cm,
    ymin=0, ymax=1.1,
    ytick={0, 0.2, 0.4, 0.6, 0.8, 1.0},
    ylabel={Recognition rate},
    xlabel={Animation},
    xtick=data,
    symbolic x coords={A1, A2, A3, A4, A5, A6},
    xticklabels={
        {Curiosity},
        {Happiness},
        {Sadness},
        {Frustrated},
        {Anger},
        {Fear/Shy}
    },
    x tick label style={
        font=\small,
        rotate=20,
        anchor=east,
        yshift=-4pt,
    },
    legend style={
        at={(0.5,-0.35)},
        anchor=north,
        legend columns=2,
        font=\small,
        column sep=10pt,
    },
    enlarge x limits=0.15,
    tick label style={font=\small},
    label style={font=\small},
    xlabel style={yshift=-10pt},
    ymajorgrids=true,
    grid style={line width=0.2pt, draw=gray!30},
    nodes near coords,
    nodes near coords align={vertical},
    every node near coord/.append style={font=\tiny},
    point meta=explicit symbolic,
]

% Live study — navy
\addplot+[ybar, fill=navyblue, draw=white, line width=0.5pt]
    coordinates {
        (A1, 0.25) [0.25]
        (A2, 1.00) [1.00]
        (A3, 0.88) [0.88]
        (A4, 0.75) [0.75]
        (A5, 1.00) [1.00]
        (A6, 0.81) [0.81]
    };

% Video study — orange
\addplot+[ybar, fill=emotion, draw=white, line width=0.5pt]
    coordinates {
        (A1, 0.05) [0.05]
        (A2, 0.95) [0.95]
        (A3, 0.95) [0.95]
        (A4, 0.32) [0.32]
        (A5, 0.95) [0.95]
        (A6, 0.36) [0.36]
    };

\legend{Live study (N=16), Video study (N=22)}

\end{axis}
\end{tikzpicture}
\caption[Emotion recognition rates per animation for the live
study]{Emotion recognition rates per animation for the live
study (N=11) and video-based study (N=22). Values represent
The percentage of participants who correctly identified the
intended emotion.}
\label{fig:recognition_rate}
\end{figure}

By examining these results, a clear distinction between the animations can be seen. Some animations were more consistently recognized, while others proved more ambiguous in interpretation. \textit{Happiness} and \textit{Anger} achieve near-perfect scores in both live and video-based studies, with live rates of 100\% and video rates of 95\%. Sadness also offers a high recognition rate, with 88\% in the live study and 95\% in the video-based study. The \textit{Curiosity} animation was the weakest performer overall, achieving recognition rates of 25\% in the live study and only 5\% in the video-based study. Results also show that the live study consistently outperformed the video-based study across all animations, except for the \textit{sadness} animation. A notable performance gap was also observed for \textit{Disappointment} (75\% live vs.\ 32\% video) and \textit{Fear} 
(81\% live vs.\ 36\% video), which are examined in detail below.

Figure ~\ref{fig:emotion_recognition_video} shows how the different responses were classified for each of the 6 emotions. \textit{Anger}, \textit{Happiness}, and \textit{Frustrated} show a low level of variation consistent with their high recognition rates. The remaining 3 animations display a wider variety of answers, indicating that participants interpreted them differently.




\begin{figure}[H]
\centering
\begin{tikzpicture}
\begin{axis}[
    ybar stacked,
    bar width=36pt,
    width=14cm,
    height=9cm,
    ymin=0, ymax=22,
    ytick={0,5,10,15,20,22},
    ylabel={Number of responses},
    xlabel={Animation},
    xtick=data,
    symbolic x coords={A1, A2, A3, A4, A5, A6},
    xticklabels={
        {Curiosity},
        {Happiness},
        {Sadness},
        {Frustrated},
        {Anger},
        {Fear}
    },
    x tick label style={
        font=\small,
        rotate=20,
        anchor=east,
        yshift=-5pt,
    },
    legend style={
        at={(0.5,-0.28)},
        anchor=north,
        legend columns=3,
        font=\small,
        column sep=10pt,
    },
    enlarge x limits=0.12,
    tick label style={font=\small},
    label style={font=\small},
    xlabel style={yshift=-10pt},
    ymajorgrids=true,
    grid style={line width=0.2pt, draw=gray!30},
]

\addplot+[ybar, fill=intended, draw=white, line width=0.5pt]
    coordinates {
        (A1, 1)(A2, 21)(A3, 21)(A4, 7)(A5, 21)(A6, 8)
    };

\addplot+[ybar, fill=adjacent, draw=white, line width=0.5pt]
    coordinates {
        (A1, 10)(A2, 0)(A3, 1)(A4, 14)(A5, 0)(A6, 6)
    };

\addplot+[ybar, fill=other, draw=white, line width=0.5pt]
    coordinates {
        (A1, 11)(A2, 1)(A3, 0)(A4, 1)(A5, 1)(A6, 8)
    };

\legend{
    Intended emotion,
    Adjacent/related emotion,
    Other
}

\end{axis}
\end{tikzpicture}
\caption[Emotion recognition results in video-based study]{Distribution of emotion recognition responses across all six
animations for the video-based study (N=22). Responses were grouped
into three categories: intended emotion (matching the target),
adjacent/related emotion (overlapping emotional state), and other
(unrelated responses). Dark purple indicates correct identification
of the intended emotion.}
\label{fig:emotion_recognition_video}
\end{figure}
\newpage
\begin{figure}[H]
\centering
\begin{tikzpicture}
\begin{axis}[
    ybar stacked,
    bar width=36pt,
    width=14cm,
    height=9cm,
    ymin=0, ymax=11,
    ytick={0,2,4,6,8,10,11},
    ylabel={Number of responses},
    xlabel={Animation},
    xtick=data,
    symbolic x coords={A1, A2, A3, A4, A5, A6},
    xticklabels={
        {Curiosity},
        {Happiness},
        {Sadness},
        {Frustrated},
        {Anger},
        {Fear}
    },
    x tick label style={
        font=\small,
        rotate=20,
        anchor=east,
        yshift=-5pt,
    },
    legend style={
        at={(0.5,-0.28)},
        anchor=north,
        legend columns=3,
        font=\small,
        column sep=10pt,
    },
    enlarge x limits=0.12,
    tick label style={font=\small},
    label style={font=\small},
    xlabel style={yshift=-10pt},
    ymajorgrids=true,
    grid style={line width=0.2pt, draw=gray!30},
]

\addplot+[ybar, fill=intended, draw=white, line width=0.5pt]
    coordinates {
        (A1, 4)(A2, 11)(A3, 10)(A4, 9)(A5, 11)(A6, 8)
    };

\addplot+[ybar, fill=adjacent, draw=white, line width=0.5pt]
    coordinates {
        (A1, 5)(A2, 0)(A3, 1)(A4, 2)(A5, 0)(A6, 2)
    };

\addplot+[ybar, fill=other, draw=white, line width=0.5pt]
    coordinates {
        (A1, 2)(A2, 0)(A3, 0)(A4, 0)(A5, 0)(A6, 1)
    };

\legend{
    Intended emotion,
    Adjacent/related emotion,
    Other
}

\end{axis}
\end{tikzpicture}
\caption{Distribution of emotion recognition responses across all six
animations for the live study (N=11). Responses were grouped
into three categories: intended emotion (matching the target),
adjacent/related emotion (overlapping emotional state), and other
(unrelated responses). Dark purple indicates correct identification
of the intended emotion.}
\label{fig:emotion_recognition_live}
\end{figure}
\textbf{Curiosity} had the highest concentration of responses categorized as \textit{other}, with recurring answers including confusion, tiredness, sadness, and frustration. This was also the weakest performing animation overall, achieving recognition rates of only 25\% in the live study and 5\% in the video-based study.

\textbf{Frustrated} had the highest concentration of \textit{adjacent} responses while still having a very low concentration of \textit{other} responses. Recurring answers included impatience, disgust, and feeling upset.

\textbf{Fear} also displayed a wide spread of responses across all three categories. In the live study, 81\% of participants correctly identified the intended emotion, whereas the video-based study achieved only 36\%, representing the largest gap between the two groups, particularly for \textit{Frustrated}. Recurring adjacent responses included shyness, sadness, and surprise.





\newpage
\section{Likert Scale Results}
% Confidence, expressiveness, naturalness per animation

\begin{figure}[H]
\centering
\begin{tikzpicture}
\begin{axis}[
    ybar,
    bar width=8pt,
    width=14cm,
    height=9cm,
    ymin=0, ymax=5,
    ytick={0,1,2,3,4,5},
    ylabel={Mean score (1--5)},
    xlabel={Animation},
    xtick=data,
    symbolic x coords={A1, A2, A3, A4, A5, A6},
    xticklabels={
        {Curiosity},
        {Happiness},
        {Sadness},
        {Frustrated},
        {Anger},
        {Fear}
    },
    x tick label style={
        font=\small,
        rotate=20,
        anchor=east,
        yshift=-4pt,
    },
    legend style={
        at={(0.5,-0.35)},
        anchor=north,
        legend columns=3,
        font=\small,
        column sep=10pt,
    },
    enlarge x limits=0.12,
    tick label style={font=\small},
    label style={font=\small},
    xlabel style={yshift=-10pt},
    ymajorgrids=true,
    grid style={line width=0.2pt, draw=gray!30},
]

% Confidence — dark purple
\addplot+[ybar, fill=intended, draw=white, line width=0.5pt]
    coordinates {
        (A1,3.09)(A2,4.23)(A3,4.45)(A4,3.18)(A5,4.50)(A6,3.23)
    };

% Expressiveness — lavender
\addplot+[ybar, fill=adjacent, draw=white, line width=0.5pt]
    coordinates {
        (A1,3.82)(A2,4.05)(A3,4.55)(A4,3.32)(A5,4.41)(A6,3.73)
    };

% Naturalness — orange
\addplot+[ybar, fill=other, draw=white, line width=0.5pt]
    coordinates {
        (A1,3.50)(A2,3.64)(A3,4.23)(A4,3.32)(A5,4.18)(A6,3.82)
    };

\legend{
    Confidence,
    Expressiveness,
    Naturalness
}

\end{axis}
\end{tikzpicture}
\caption{Mean Likert scale scores per animation for the video-based 
study (N=22). Each animation shows three bars representing mean 
confidence, expressiveness, and naturalness ratings on a scale of 
1 to 5.}
\label{fig:likert_video}
\end{figure}
\begin{figure}[H]
\centering
\begin{tikzpicture}
\begin{axis}[
    ybar,
    bar width=8pt,
    width=14cm,
    height=9cm,
    ymin=0, ymax=5,
    ytick={0,1,2,3,4,5},
    ylabel={Mean score (1--5)},
    xlabel={Animation},
    xtick=data,
    symbolic x coords={A1, A2, A3, A4, A5, A6},
    xticklabels={
        {Curiosity},
        {Happiness},
        {Sadness},
        {Frustrated},
        {Anger},
        {Fear}
    },
    x tick label style={
        font=\small,
        rotate=20,
        anchor=east,
        yshift=-4pt,
    },
    legend style={
        at={(0.5,-0.35)},
        anchor=north,
        legend columns=3,
        font=\small,
        column sep=10pt,
    },
    enlarge x limits=0.12,
    tick label style={font=\small},
    label style={font=\small},
    xlabel style={yshift=-10pt},
    ymajorgrids=true,
    grid style={line width=0.2pt, draw=gray!30},
]

% Confidence — dark purple
\addplot+[ybar, fill=intended, draw=white, line width=0.5pt]
    coordinates {
        (A1,3.55)(A2,4.18)(A3,4.64)(A4,4.09)(A5,4.73)(A6,2.82)
    };

% Expressiveness — lavender
\addplot+[ybar, fill=adjacent, draw=white, line width=0.5pt]
    coordinates {
        (A1,3.82)(A2,4.27)(A3,4.55)(A4,4.09)(A5,4.36)(A6,3.36)
    };

% Naturalness — orange
\addplot+[ybar, fill=other, draw=white, line width=0.5pt]
    coordinates {
        (A1,3.73)(A2,4.09)(A3,4.18)(A4,4.18)(A5,4.45)(A6,3.45)
    };

\legend{
    Confidence,
    Expressiveness,
    Naturalness
}

\end{axis}
\end{tikzpicture}
\caption[Mean Likert scale scores per animation for the live 
study]{Mean Likert scale scores per animation for the live 
study (N=11). Each animation shows three bars representing mean 
confidence, expressiveness, and naturalness ratings on a scale of 
1 to 5.}
\label{fig:likert_live}
\end{figure}




\subsection*{Godspeed Results}
% Five dimensions ( spider chart)
\begin{figure}[H]
\centering
\begin{tikzpicture}[scale=3.0]
\def\N{5}
\def\levels{5}
% Colors
\colorlet{spiderline}{emotion}
\colorlet{spiderfill}{emotion}
% Axis labels and values
\def\labels{{"Anthropomorphism", "Animacy", "Likeability", 
             "Perceived\\Intelligence", "Perceived\\Safety"}}
\def\values{{3.15, 3.49, 4.45, 3.45, 4.19}}
% Draw background grid rings
\foreach \lev in {1,...,5}{
    \draw[gray!25, line width=0.4pt]
        \foreach \i in {1,...,5}{
            -- ({(\lev/5)*cos(90+(\i-1)*360/5)},
               {(\lev/5)*sin(90+(\i-1)*360/5)})
        } -- cycle;
}
% Draw axis lines
\foreach \i in {1,...,5}{
    \draw[gray!35, line width=0.4pt]
        (0,0) -- ({cos(90+(\i-1)*360/5)}, {sin(90+(\i-1)*360/5)});
}
% Draw tick labels on one axis (right side)
\foreach \lev in {1,...,5}{
    \node[font=\tiny, gray] at 
        ({(\lev/5)*cos(90)}, {(\lev/5)*sin(90)+0.06}) {\lev};
}
% Plot data — video study
\fill[spiderfill, opacity=0.2]
    ({(3.15/5)*cos(90)},   {(3.15/5)*sin(90)})   --
    ({(3.49/5)*cos(90-72)}, {(3.49/5)*sin(90-72)}) --
    ({(4.45/5)*cos(90-144)},{(4.45/5)*sin(90-144)}) --
    ({(3.45/5)*cos(90-216)},{(3.45/5)*sin(90-216)}) --
    ({(4.19/5)*cos(90-288)},{(4.19/5)*sin(90-288)}) --
    cycle;
\draw[spiderline, line width=1.8pt]
    ({(3.15/5)*cos(90)},   {(3.15/5)*sin(90)})   --
    ({(3.49/5)*cos(90-72)}, {(3.49/5)*sin(90-72)}) --
    ({(4.45/5)*cos(90-144)},{(4.45/5)*sin(90-144)}) --
    ({(3.45/5)*cos(90-216)},{(3.45/5)*sin(90-216)}) --
    ({(4.19/5)*cos(90-288)},{(4.19/5)*sin(90-288)}) --
    cycle;
% Data point dots
\foreach \val/\angle in {
    3.15/90, 3.49/18, 4.45/-54, 3.45/-126, 4.19/-198}{
    \filldraw[spiderline] 
        ({(\val/5)*cos(\angle)}, {(\val/5)*sin(\angle)}) 
        circle (1pt);
}
% Value labels next to dots
\node[font=\tiny, above] at 
    ({(3.15/5)*cos(90)+0.0},   {(3.15/5)*sin(90)+0.08})   {3.15};
\node[font=\tiny, right] at 
    ({(3.49/5)*cos(18)+0.06},  {(3.49/5)*sin(18)})        {3.49};
\node[font=\tiny, right] at 
    ({(4.45/5)*cos(-54)+0.06}, {(4.45/5)*sin(-54)})       {4.45};
\node[font=\tiny, left] at 
    ({(3.45/5)*cos(-126)-0.06},{(3.45/5)*sin(-126)})      {3.45};
\node[font=\tiny, left] at 
    ({(4.19/5)*cos(-198)-0.06},{(4.19/5)*sin(-198)})      {4.19};
% Axis labels
\node[font=\small, above, align=center] at 
    ({1.22*cos(90)},   {1.22*sin(90)})    {Anthropomorphism};
\node[font=\small, right, align=left] at 
    ({1.22*cos(18)},   {1.22*sin(18)})    {Animacy};
\node[font=\small, right, align=left] at 
    ({1.22*cos(-54)},  {1.22*sin(-54)})   {Likeability};
\node[font=\small, left, align=right] at 
    ({1.22*cos(-126)}, {1.22*sin(-126)})  {Perceived\\Intelligence};
\node[font=\small, left, align=right] at 
    ({1.22*cos(-198)}, {1.22*sin(-198)})  {Perceived\\Safety};
\end{tikzpicture}
\caption{Godspeed questionnaire results for the video-based study 
(N=22) across five dimensions. Scores are on a scale of 1 to 5.}
\label{fig:godspeed_video}
\end{figure} 

\begin{figure}[H]
\centering
\begin{tikzpicture}[scale=3.0]
\def\N{5}
\def\levels{5}
% Colors
\colorlet{spiderline}{navyblue}
\colorlet{spiderfill}{navyblue}
% Draw background grid rings
\foreach \lev in {1,...,5}{
    \draw[gray!25, line width=0.4pt]
        \foreach \i in {1,...,5}{
            -- ({(\lev/5)*cos(90+(\i-1)*360/5)},
               {(\lev/5)*sin(90+(\i-1)*360/5)})
        } -- cycle;
}
% Draw axis lines
\foreach \i in {1,...,5}{
    \draw[gray!35, line width=0.4pt]
        (0,0) -- ({cos(90+(\i-1)*360/5)}, {sin(90+(\i-1)*360/5)});
}
% Draw tick labels
\foreach \lev in {1,...,5}{
    \node[font=\tiny, gray] at 
        ({(\lev/5)*cos(90)}, {(\lev/5)*sin(90)+0.06}) {\lev};
}
% Plot data — live study N=16
\fill[spiderfill, opacity=0.2]
    ({(3.14/5)*cos(90)},   {(3.14/5)*sin(90)})   --
    ({(3.40/5)*cos(90-72)}, {(3.40/5)*sin(90-72)}) --
    ({(4.50/5)*cos(90-144)},{(4.50/5)*sin(90-144)}) --
    ({(3.20/5)*cos(90-216)},{(3.20/5)*sin(90-216)}) --
    ({(4.26/5)*cos(90-288)},{(4.26/5)*sin(90-288)}) --
    cycle;
\draw[spiderline, line width=1.8pt]
    ({(3.14/5)*cos(90)},   {(3.14/5)*sin(90)})   --
    ({(3.40/5)*cos(90-72)}, {(3.40/5)*sin(90-72)}) --
    ({(4.50/5)*cos(90-144)},{(4.50/5)*sin(90-144)}) --
    ({(3.20/5)*cos(90-216)},{(3.20/5)*sin(90-216)}) --
    ({(4.26/5)*cos(90-288)},{(4.26/5)*sin(90-288)}) --
    cycle;
% Data point dots
\foreach \val/\angle in {
    3.14/90, 3.40/18, 4.50/-54, 3.20/-126, 4.26/-198}{
    \filldraw[spiderline] 
        ({(\val/5)*cos(\angle)}, {(\val/5)*sin(\angle)}) 
        circle (1pt);
}
% Value labels
\node[font=\tiny, above] at 
    ({(3.14/5)*cos(90)+0.0},   {(3.14/5)*sin(90)+0.08})   {3.14};
\node[font=\tiny, right] at 
    ({(3.40/5)*cos(18)+0.06},  {(3.40/5)*sin(18)})        {3.40};
\node[font=\tiny, right] at 
    ({(4.50/5)*cos(-54)+0.06}, {(4.50/5)*sin(-54)})       {4.50};
\node[font=\tiny, left] at 
    ({(3.20/5)*cos(-126)-0.06},{(3.20/5)*sin(-126)})      {3.20};
\node[font=\tiny, left] at 
    ({(4.26/5)*cos(-198)-0.06},{(4.26/5)*sin(-198)})      {4.26};
% Axis labels
\node[font=\small, above, align=center] at 
    ({1.22*cos(90)},   {1.22*sin(90)})    {Anthropomorphism};
\node[font=\small, right, align=left] at 
    ({1.22*cos(18)},   {1.22*sin(18)})    {Animacy};
\node[font=\small, right, align=left] at 
    ({1.22*cos(-54)},  {1.22*sin(-54)})   {Likeability};
\node[font=\small, left, align=right] at 
    ({1.22*cos(-126)}, {1.22*sin(-126)})  {Perceived\\Intelligence};
\node[font=\small, left, align=right] at 
    ({1.22*cos(-198)}, {1.22*sin(-198)})  {Perceived\\Safety};
\end{tikzpicture}
\caption[Godspeed questionnaire results for the live-based study]{Godspeed 
questionnaire results for the live study (N=16) across five dimensions. 
Scores are on a scale of 1 to 5.}
\label{fig:godspeed_live}
\end{figure}

\begin{figure}[H]
\centering
\begin{tikzpicture}[scale=3.0]

\foreach \lev in {1,...,5}{
    \draw[gray!25, line width=0.4pt]
        \foreach \i in {1,...,5}{
            -- ({(\lev/5)*cos(90+(\i-1)*360/5)},
               {(\lev/5)*sin(90+(\i-1)*360/5)})
        } -- cycle;
}
\foreach \i in {1,...,5}{
    \draw[gray!35, line width=0.4pt]
        (0,0) -- ({cos(90+(\i-1)*360/5)}, {sin(90+(\i-1)*360/5)});
}
\foreach \lev in {1,...,5}{
    \node[font=\tiny, gray] at
        ({(\lev/5)*cos(90)}, {(\lev/5)*sin(90)+0.06}) {\lev};
}

% Live — navy solid
\fill[navyblue, opacity=0.15]
    ({(3.00/5)*cos(90)},   {(3.00/5)*sin(90)})   --
    ({(3.38/5)*cos(18)},   {(3.38/5)*sin(18)})   --
    ({(4.58/5)*cos(-54)},  {(4.58/5)*sin(-54)})  --
    ({(3.20/5)*cos(-126)}, {(3.20/5)*sin(-126)}) --
    ({(4.24/5)*cos(-198)}, {(4.24/5)*sin(-198)}) --
    cycle;
\draw[navyblue, line width=1.8pt]
    ({(3.00/5)*cos(90)},   {(3.00/5)*sin(90)})   --
    ({(3.38/5)*cos(18)},   {(3.38/5)*sin(18)})   --
    ({(4.58/5)*cos(-54)},  {(4.58/5)*sin(-54)})  --
    ({(3.20/5)*cos(-126)}, {(3.20/5)*sin(-126)}) --
    ({(4.24/5)*cos(-198)}, {(4.24/5)*sin(-198)}) --
    cycle;
\foreach \val/\angle in {
    3.00/90, 3.38/18, 4.58/-54, 3.20/-126, 4.24/-198}{
    \filldraw[navyblue]
        ({(\val/5)*cos(\angle)}, {(\val/5)*sin(\angle)}) circle (1.0pt);
}

% Video — orange dashed
\fill[emotion, opacity=0.15]
    ({(3.15/5)*cos(90)},   {(3.15/5)*sin(90)})   --
    ({(3.49/5)*cos(18)},   {(3.49/5)*sin(18)})   --
    ({(4.45/5)*cos(-54)},  {(4.45/5)*sin(-54)})  --
    ({(3.45/5)*cos(-126)}, {(3.45/5)*sin(-126)}) --
    ({(4.19/5)*cos(-198)}, {(4.19/5)*sin(-198)}) --
    cycle;
\draw[emotion, line width=1.8pt, dashed]
    ({(3.15/5)*cos(90)},   {(3.15/5)*sin(90)})   --
    ({(3.49/5)*cos(18)},   {(3.49/5)*sin(18)})   --
    ({(4.45/5)*cos(-54)},  {(4.45/5)*sin(-54)})  --
    ({(3.45/5)*cos(-126)}, {(3.45/5)*sin(-126)}) --
    ({(4.19/5)*cos(-198)}, {(4.19/5)*sin(-198)}) --
    cycle;
\foreach \val/\angle in {
    3.15/90, 3.49/18, 4.45/-54, 3.45/-126, 4.19/-198}{
    \filldraw[emotion]
        ({(\val/5)*cos(\angle)}, {(\val/5)*sin(\angle)}) circle (1.0pt);
}

% Axis labels
\node[font=\small, above, align=center] at ({1.28*cos(90)},{1.28*sin(90)})   {Anthropomorphism};
\node[font=\small, right, align=left]   at ({1.28*cos(18)},{1.28*sin(18)})   {Animacy};
\node[font=\small, right, align=left]   at ({1.28*cos(-54)},{1.28*sin(-54)}) {Likeability};
\node[font=\small, left, align=right]   at ({1.28*cos(-126)},{1.28*sin(-126)}) {Perceived\\Intelligence};
\node[font=\small, left, align=right]   at ({1.28*cos(-198)},{1.28*sin(-198)}) {Perceived\\Safety};

% Legend
\draw[navyblue, line width=1.5pt] (-1.1,-1.25) -- (-0.7,-1.25);
\filldraw[navyblue] (-0.9,-1.25) circle (1.2pt);
\node[font=\small, right] at (-0.68,-1.25) {Live study (N=11)};

\draw[emotion, line width=1.5pt, dashed] (-1.1,-1.45) -- (-0.7,-1.45);
\filldraw[emotion] (-0.9,-1.45) circle (1.0pt);
\node[font=\small, right] at (-0.68,-1.45) {Video study (N=22)};

\end{tikzpicture}
\caption{Godspeed questionnaire results for both study groups across
five dimensions. Solid line represents the live study (N=11) and
dashed line represents the video-based study (N=22).
Scores are on a scale of 1 to 5.}
\label{fig:godspeed_overlay}
\end{figure}





%\section{RL Locomotion Experiments}
%\subsection{Simulation experiments}
%\subsubsection{Gait performance in simulation}
%\subsubsection{Hardware gait performance}